\documentclass[11pt,a4paper]{article}
\usepackage{get-uebung}

\blatt{3}{Kirchhoffsche Sätze}{Kap.~2.1}

\begin{document}

\blattkopf{
  \item den Knotensatz an einzelnen Knoten und an Hüllflächen vorzeichenrichtig anwenden,
  \item mit dem Maschensatz unbekannte Spannungen und Knotenpotentiale bestimmen,
  \item für ein Netzwerk mit zwei Quellen die nötige Anzahl unabhängiger Knoten- und Maschengleichungen angeben und das Gleichungssystem lösen,
  \item negative Ergebnisse richtig interpretieren.
}

% -------------------------------------------------------------------
\begin{aufgabe}{Knotensatz}{12}
\begin{teilaufgaben}
\item An einem Knoten treffen fünf Leitungen zusammen (Bild links). Gegeben sind $I_1 = \SI{3}{A}$, $I_2 = \SI{-1}{A}$, $I_3 = \SI{2}{A}$ und $I_5 = \SI{0.5}{A}$. Berechnen Sie $I_4$.
\item Im Netzwerk rechts sind $I_\text{ein} = \SI{6}{A}$, $I_a = \SI{4}{A}$ und $I_d = \SI{1.5}{A}$ bekannt. Bestimmen Sie $I_b$, $I_c$, $I_e$ und $I_\text{aus}$. Die Werte der Widerstände werden nicht benötigt.
\item Kontrollieren Sie Ihr Ergebnis aus b) mit einer Hüllfläche um das gesamte Netzwerk.
\end{teilaufgaben}
\begin{center}
\begin{tikzpicture}[>=Latex,baseline=(k.center)]
  \node[circle,fill,inner sep=1.8pt] (k) at (0,0) {};
  \draw[thick,->] (-2,0) node[left]{$I_1$} -- (-0.12,0);
  \draw[thick,->] (-1.3,1.3) node[above left]{$I_2$} -- (-0.1,0.1);
  \draw[thick,->] (0.1,0.1) -- (1.3,1.3) node[above right]{$I_3$};
  \draw[thick,->] (0.12,0) -- (2,0) node[right]{$I_4$};
  \draw[thick,->] (0,-1.6) node[below]{$I_5$} -- (0,-0.12);
\end{tikzpicture}
\hspace{1.2cm}
\begin{circuitikz}[baseline=(c.center)]
  \coordinate (c) at (2,1.5);
  \draw (-2.4,3) to[short, i=$I_\text{ein}$, o-] (0,3)
        to[R, l=$R_1$, i=$I_a$, *-*] (4,3);
  \draw (0,3) to[R, l_=$R_2$, i=$I_b$] (0,0);
  \draw (4,0) to[R, l=$R_3$, i_=$I_c$, *-] (4,3);
  \draw (0,0) to[R, l^=$R_4$, i_=$I_d$, *-] (4,0);
  \draw (4,0) to[short, i=$I_\text{aus}$, -o] (6.4,0);
  \draw (0,0) to[short, i=$I_e$, -o] (-2.4,0);
  \node[below right] at (0,3) {$K_1$};
  \node[below left] at (4,3) {$K_2$};
  \node[above left] at (4,0) {$K_3$};
  \node[above left] at (0,0) {$K_4$};
\end{circuitikz}
\end{center}
\begin{tipp}Zählen Sie Ströme, deren Pfeil \emph{in} den Knoten zeigt, positiv und solche, deren Pfeil herauszeigt, negativ -- unabhängig vom Vorzeichen ihres Zahlenwerts.\end{tipp}
\end{aufgabe}

\begin{loesung}
\schritt{a) Einzelner Knoten}
\leitfrage{Welche Pfeile zeigen in den Knoten hinein, welche heraus?}
Hinein: $I_1$, $I_2$, $I_5$. Heraus: $I_3$, $I_4$. Knotensatz:
\[
  I_1 + I_2 + I_5 - I_3 - I_4 = 0
  \quad\Rightarrow\quad
  I_4 = I_1 + I_2 + I_5 - I_3 = 3 + (-1) + 0{,}5 - 2\ \si{A} = \ergebnis{\SI{0.5}{A}}
\]
Den negativen Wert von $I_2$ setzt man einfach \emph{mit} seinem Vorzeichen ein: In Wirklichkeit fließt dort $\SI{1}{A}$ aus dem Knoten heraus.
\schritt{b) Netzwerk -- Knoten für Knoten}
\leitfrage{Bei welchem Knoten ist nur \emph{ein} Strom unbekannt? Dort beginnen.}
Knoten $K_1$: $I_\text{ein} - I_a - I_b = 0 \;\Rightarrow\; I_b = 6 - 4\ \si{A} = \ergebnis{\SI{2}{A}}$\\[2pt]
Knoten $K_2$: Beide Pfeile ($I_a$ und $I_c$) zeigen hinein: $I_a + I_c = 0 \;\Rightarrow\; \ergebnis{I_c = \SI{-4}{A}}$\\
\phantom{Knoten $K_2$:} Der Strom fließt in $R_3$ also tatsächlich von oben nach unten.\\[2pt]
Knoten $K_4$: $I_b - I_d - I_e = 0 \;\Rightarrow\; I_e = 2 - 1{,}5\ \si{A} = \ergebnis{\SI{0.5}{A}}$\\[2pt]
Knoten $K_3$: hinein $I_d$, heraus $I_c$ und $I_\text{aus}$:
\[
  I_d - I_c - I_\text{aus} = 0 \quad\Rightarrow\quad I_\text{aus} = 1{,}5 - (-4)\ \si{A} = \ergebnis{\SI{5.5}{A}}
\]
\schritt{c) Hüllfläche}
Der Knotensatz gilt auch für eine geschlossene Fläche um einen ganzen Schaltungsteil. Sie wird nur von $I_\text{ein}$, $I_\text{aus}$ und $I_e$ durchstoßen:
\[
  I_\text{ein} - I_\text{aus} - I_e = 6 - 5{,}5 - 0{,}5\ \si{A} = 0 \;\checkmark
\]
\typfehler{Bei $K_3$ das Minuszeichen „doppelt“ vergessen: $I_c$ zeigt heraus \emph{und} ist negativ, also $-I_c = +\SI{4}{A}$. Das Vorzeichen in der Gleichung kommt vom \emph{Pfeil}, das Vorzeichen des Werts kommt beim Einsetzen hinzu.}
\end{loesung}

% -------------------------------------------------------------------
\begin{aufgabe}{Maschensatz und Potentiale}{13}
In der Schaltung sind $U_0 = \SI{12}{V}$, $U_1 = \SI{5}{V}$ und $U_3 = \SI{-3}{V}$ bekannt. Die Zweipole $Z_1$ bis $Z_4$ sind nicht näher bekannt.
\begin{center}
\begin{circuitikz}
  \draw (0,0) to[V, v<=$U_0$] (0,3)
        to[R, l=$Z_1$, v=$U_1$, -*] (3,3)
        to[R, l=$Z_3$, v<=$U_3$, -*] (6,3)
        to[R, l=$Z_4$, v=$U_4$] (6,0) -- (0,0);
  \draw (3,3) to[R, l=$Z_2$, v=$U_2$, -*] (3,0);
  \draw (3,0) -- (3,-0.3) node[ground]{};
  \node[above left] at (0,3) {A};
  \node[above] at (3,3.05) {B};
  \node[above right] at (6,3) {C};
  \node[below left] at (3,-0.05) {0};
\end{circuitikz}
\end{center}
\begin{teilaufgaben}
  \item Zeichnen Sie zwei Maschen mit Umlaufrichtung ein und berechnen Sie $U_2$ und $U_4$.
  \item Kontrollieren Sie das Ergebnis mit der äußeren Masche.
  \item Der untere Knoten 0 ist der Bezugsknoten ($\varphi_0 = \SI{0}{V}$). Bestimmen Sie die Potentiale $\varphi_A$, $\varphi_B$ und $\varphi_C$.
  \item Welche Spannung $U_{AC}$ zeigt ein Spannungsmesser zwischen A und C an?
\end{teilaufgaben}
\begin{tipp}Ein Spannungspfeil von Punkt X nach Punkt Y bedeutet $U = \varphi_X - \varphi_Y$.\end{tipp}
\end{aufgabe}

\begin{loesung}
\schritt{a) Zwei Maschen im Uhrzeigersinn}
\leitfrage{Laufen Sie die linke Masche im Uhrzeigersinn ab (links hoch, oben nach rechts, durch $Z_2$ hinunter). Welche Pfeile zeigen in Umlaufrichtung?}
Masche I (A--B--0--A): $U_0$ zeigt nach unten, also \emph{gegen} die Umlaufrichtung; $U_1$ und $U_2$ zeigen mit ihr.
\[
  -U_0 + U_1 + U_2 = 0 \quad\Rightarrow\quad U_2 = U_0 - U_1 = 12 - 5\ \si{V} = \ergebnis{\SI{7}{V}}
\]
Masche II (B--C--0--B): durch $Z_2$ hinauf (gegen $U_2$), oben nach rechts (gegen $U_3$, der nach links zeigt), durch $Z_4$ hinunter (mit $U_4$):
\[
  -U_2 - U_3 + U_4 = 0 \quad\Rightarrow\quad U_4 = U_2 + U_3 = 7 + (-3)\ \si{V} = \ergebnis{\SI{4}{V}}
\]
\schritt{b) Kontrolle: äußere Masche}
A--B--C--0--A: $-U_0 + U_1 - U_3 + U_4 = -12 + 5 - (-3) + 4\ \si{V} = 0$ \checkmark
\schritt{c) Potentiale}
Wir gehen vom Bezugsknoten aus und nutzen $U = \varphi_X - \varphi_Y$:
\begin{align*}
  U_0 = \varphi_A - \varphi_0 &\Rightarrow \ergebnis{\varphi_A = \SI{12}{V}}\\
  U_2 = \varphi_B - \varphi_0 &\Rightarrow \ergebnis{\varphi_B = \SI{7}{V}}\\
  U_4 = \varphi_C - \varphi_0 &\Rightarrow \ergebnis{\varphi_C = \SI{4}{V}}
\end{align*}
Probe mit den übrigen Spannungen: $U_1 = \varphi_A - \varphi_B = \SI{5}{V}$ \checkmark, $U_3 = \varphi_C - \varphi_B = \SI{-3}{V}$ \checkmark
\schritt{d) Spannung zwischen A und C}
\[
  U_{AC} = \varphi_A - \varphi_C = 12 - 4\ \si{V} = \ergebnis{\SI{8}{V}}
\]
Mit Potentialen braucht man dafür keinen Maschenumlauf mehr -- jede Spannung ist einfach eine Potentialdifferenz.
\typfehler{Bei Masche II $+U_3$ statt $-U_3$ schreiben, weil man nur auf das Vorzeichen des \emph{Werts} achtet. Erst die Gleichung mit Pfeilrichtungen aufstellen, \emph{dann} Zahlen mit Vorzeichen einsetzen.}
\end{loesung}

% -------------------------------------------------------------------
\begin{aufgabe}{Netzwerk mit zwei Quellen}{25}
Ein Ladegerät ($U_1 = \SI{20}{V}$, Innenwiderstand $R_1 = \SI{1}{\ohm}$) und ein Akku ($U_2 = \SI{6}{V}$, Innenwiderstand $R_2 = \SI{3}{\ohm}$) versorgen gemeinsam einen Verbraucher $R_3 = \SI{2}{\ohm}$.
\begin{center}
\begin{circuitikz}
  \draw (0,0) to[V, v<=$U_1$] (0,3)
        to[R, l=$R_1$, v_=$U_{R1}$] (2.6,3) to[short, i=$I_1$] (4,3);
  \draw (8,0) to[V, v_<=$U_2$] (8,3)
        to[R, l_=$R_2$, v^=$U_{R2}$] (5.4,3) to[short, i=$I_2$] (4,3);
  \draw (4,3) to[R, l=$R_3$, v=$U_{R3}$, *-] (4,0.9) to[short, i=$I_3$, -*] (4,0);
  \draw (0,0) -- (8,0);
\end{circuitikz}
\end{center}
\begin{teilaufgaben}
  \item Wie viele Knoten $k$ und Zweige $z$ hat die Schaltung? Wie viele unabhängige Knoten- und Maschengleichungen benötigen Sie, um alle Zweigströme zu bestimmen?
  \item Stellen Sie die Knotengleichung und die Maschengleichungen auf. Verwenden Sie dabei das Ohmsche Gesetz, sodass nur noch $I_1$, $I_2$, $I_3$ als Unbekannte auftreten.
  \item Lösen Sie das Gleichungssystem.
  \item Interpretieren Sie das Vorzeichen von $I_2$. Was geschieht mit dem Akku?
  \item Stellen Sie eine Leistungsbilanz für alle fünf Elemente auf.
\end{teilaufgaben}
\end{aufgabe}

\begin{loesung}
\schritt{a) Knoten, Zweige, Gleichungen}
Knoten: oben Mitte und unten -- also $k = 2$ (Verbindungsleitungen gehören zum Knoten). Zweige: links, Mitte, rechts -- also $z = 3$ mit den drei Unbekannten $I_1$, $I_2$, $I_3$.
\[
  \text{Knotengleichungen: } k - 1 = \ergebnis{1}, \qquad \text{Maschengleichungen: } z - (k-1) = \ergebnis{2}
\]
Die zweite Knotengleichung (unterer Knoten) wäre nur die erste mit umgekehrtem Vorzeichen -- sie liefert nichts Neues.
\schritt{b) Gleichungen aufstellen}
Oberer Knoten ($I_1$, $I_2$ hinein, $I_3$ heraus):
\[
  \text{(K)}\quad I_1 + I_2 - I_3 = 0
\]
\leitfrage{Linke Masche im Uhrzeigersinn: $U_1$ zeigt nach unten (gegen den Umlauf), $U_{R1}$ und $U_{R3}$ zeigen mit dem Umlauf.}
\[
  \text{(M1)}\quad -U_1 + R_1 I_1 + R_3 I_3 = 0
\]
Rechte Masche im Uhrzeigersinn (in $R_3$ hinauf, nach rechts durch $R_2$, in der Quelle hinunter): $U_{R3}$ und $U_{R2}$ zeigen gegen, $U_2$ mit dem Umlauf:
\[
  \text{(M2)}\quad -R_3 I_3 - R_2 I_2 + U_2 = 0
\]
\schritt{c) Lösen}
Zahlen einsetzen und (K) als $I_3 = I_1 + I_2$ in (M1), (M2) einsetzen:
\begin{align*}
  \text{(M1)}:\quad & 1\cdot I_1 + 2\,(I_1 + I_2) = 20 \;\Rightarrow\; 3I_1 + 2I_2 = 20\\
  \text{(M2)}:\quad & 3\cdot I_2 + 2\,(I_1 + I_2) = 6 \;\Rightarrow\; 2I_1 + 5I_2 = 6
\end{align*}
(Ströme in A.) Erste Gleichung $\cdot 2$ minus zweite Gleichung $\cdot 3$: $4I_2 - 15 I_2 = 40 - 18 \Rightarrow -11 I_2 = 22$.
\[
  \ergebnis{I_2 = \SI{-2}{A}}, \qquad I_1 = \frac{20 - 2\cdot(-2)}{3}\ \si{A} = \ergebnis{\SI{8}{A}}, \qquad I_3 = I_1 + I_2 = \ergebnis{\SI{6}{A}}
\]
Probe (äußere Masche): $-U_1 + R_1 I_1 - R_2 I_2 + U_2 = -20 + 8 + 6 + 6 = 0$ \checkmark
\schritt{d) Interpretation}
$I_2 < 0$: Der Strom fließt entgegen dem Pfeil, also im Akku von oben ($+$) nach unten. Der Strom fließt in den Pluspol hinein -- der Akku wird mit $\SI{2}{A}$ \emph{geladen}.
\schritt{e) Leistungsbilanz}
Quellen (EZS, $P > 0$ heißt Abgabe): $P_{U1} = U_1 I_1 = \SI{160}{W}$ (gibt ab), $P_{U2} = U_2 I_2 = \SI{-12}{W}$ (nimmt $\SI{12}{W}$ auf).\\
Widerstände (VZS): $P_{R1} = I_1^2 R_1 = \SI{64}{W}$, $P_{R2} = I_2^2 R_2 = \SI{12}{W}$, $P_{R3} = I_3^2 R_3 = \SI{72}{W}$.
\[
  \underbrace{\SI{160}{W}}_\text{abgegeben} = \underbrace{64 + 12 + 72 + 12\ \si{W}}_\text{aufgenommen} = \SI{160}{W} \;\checkmark
\]
\typfehler{Beide Knotengleichungen verwenden und dann eine Maschengleichung „vermissen“. Bei einem negativen Ergebnis den Pfeil nachträglich umdrehen und weiterrechnen, ohne das Vorzeichen zu ändern -- lassen Sie die Pfeile stehen und rechnen Sie mit dem negativen Wert weiter.}
\end{loesung}

% -------------------------------------------------------------------
\begin{aufgabe}{Richtig oder falsch?}{7}
Entscheiden Sie und begründen Sie in einem Satz.
\begin{teilaufgaben}
  \item Ein Netzwerk mit $k$ Knoten liefert $k$ unabhängige Knotengleichungen.
  \item Die Spannung zwischen zwei Punkten hängt davon ab, über welchen Weg (welche Masche) man sie berechnet.
  \item Das Vorzeichen eines berechneten Stroms hängt von der Wahl des Zählpfeils ab, der physikalische Strom aber nicht.
  \item Wählt man einen anderen Knoten als Bezugsknoten, ändern sich alle Potentiale, aber keine einzige Spannung.
\end{teilaufgaben}
\end{aufgabe}

\begin{loesung}
\schritt{Antworten}
\begin{teilaufgaben}
  \item \textbf{Falsch.} Nur $k-1$ Gleichungen sind unabhängig; die letzte folgt aus der Summe aller anderen.
  \item \textbf{Falsch.} Genau das sagt der Maschensatz: Die Summe über jeden geschlossenen Umlauf ist null, also liefern alle Wege dieselbe Spannung.
  \item \textbf{Richtig.} Ein umgedrehter Pfeil kehrt nur das Vorzeichen des Zahlenwerts um.
  \item \textbf{Richtig.} Alle Potentiale verschieben sich um denselben Betrag; Spannungen sind Potential\emph{differenzen}, der Betrag hebt sich heraus.
\end{teilaufgaben}
\end{loesung}

\end{document}
